% Part: normal-modal-logic
% Chapter: frame-definability
% Section: equivalence-S5

\documentclass[../../../include/open-logic-section]{subfiles}

\begin{document}

\olfileid{nml}{frd}{es5}

\olsection{علاقات التكافؤ و\Log{S5}}

يوصف المنطق الموجهي~$\Log{S5}$ بأنه مجموعة ال!!{formula}s الصالحة
على جميع الأطر الكلية، أي حيث يمكن الوصول من كل عالم إلى كل عالم،
بما في ذلك نفسه. وفي هذا الوضع، يقابل~$\Box$ الضرورة ويقابل
$\Diamond$ الإمكانية: تكون~$\Box !A$ صادقة إذا كانت~$!A$ صادقة في
\emph{كل} عالم، وتكون~$\Diamond !A$ صادقة إذا كانت~$!A$ صادقة في
\emph{بعض} العوالم. ويتبين أن~$\Log{S5}$ يمكن وصفه أيضًا بأنه
ال!!{formula}s الصالحة على جميع الأطر الانعكاسية والمتناظرة والمتعدية،
أي على جميع \emph{علاقات التكافؤ}.

\begin{defn}
  تكون العلاقة الثنائية~$R$ على~$W$ \emph{علاقة تكافؤ} إذا وفقط إذا
  كانت انعكاسية ومتناظرة ومتعدية. وتكون العلاقة~$R$ على~$W$
  \emph{كلية} إذا وفقط إذا كان~$Ruv$ لكل~$u,v \in W$.
\end{defn}

لما كانت~\Ax{T} و\Ax{B} و\Ax{4} تميز الأطر الانعكاسية والمتناظرة
والمتعدية، فإن الأطر التي تكون علاقة الوصول فيها علاقة تكافؤ هي
بالضبط الأطر التي تصلح فيها ال!!{formula}s الثلاث جميعًا. ويتبين أنه
يمكن وصف علاقات التكافؤ أيضًا بتركيبات أخرى من ال!!{formula}s، لأن
الشروط التي عرّفنا بها علاقات التكافؤ تكافئ تركيبات من شروط أخرى
مألوفة على~$R$.

\begin{prop}\ollabel{prop:equivalences}
  الشروط الآتية متكافئة:
  \begin{enumerate}
  \item $R$ علاقة تكافؤ؛
  \item $R$ انعكاسية وإقليدية؛
  \item $R$ تسلسلية ومتناظرة وإقليدية؛
  \item $R$ تسلسلية ومتناظرة ومتعدية.
  \end{enumerate}
\end{prop}

\begin{proof}
  تمرين.
\end{proof}

\begin{prob}
  أثبت~\olref[nml][frd][es5]{prop:equivalences} بإثبات الآتي:
  \begin{enumerate}
  \item إذا كانت~$R$ متناظرة ومتعدية، فهي إقليدية.
  \item إذا كانت~$R$ انعكاسية، فهي تسلسلية.
  \item إذا كانت~$R$ انعكاسية وإقليدية، فهي متناظرة.
  \item إذا كانت~$R$ متناظرة وإقليدية، فهي متعدية.
  \item إذا كانت~$R$ تسلسلية ومتناظرة ومتعدية، فهي انعكاسية.
  \end{enumerate}
  اشرح لماذا يكفي هذا لإثبات تكافؤ الشروط.
\end{prob}

تمثل~\olref{prop:equivalences} النظير الدلالي لـ
\olref[prf][prs]{prop:S5}، إذ تعطي وصفًا مكافئًا للمنطق الموجهي
للأطر التي تكون~$R$ عليها علاقة تكافؤ (وهو المنطق المسمى تقليديًا
$\Log{S5}$).

ما العلاقة بين العلاقات الكلية وعلاقات التكافؤ؟ مع أن كل علاقة كلية
علاقة تكافؤ، فمن الواضح أن ليست كل علاقة تكافؤ كلية. غير أن
الصيغ الصالحة على جميع العلاقات الكلية هي نفسها بالضبط
ال!!{formula}s الصالحة على جميع علاقات التكافؤ.

\begin{prop}
  لتكن~$R$ علاقة تكافؤ، وعرّف، لكل~$w \in W$، \emph{صنف تكافؤ}~$w$
  بأنه المجموعة~$[w] = \{w'\in W : Rww'\}$. عندئذ:
  \begin{enumerate}
  \item $w \in [w]$؛
  \item $R$ كلية على كل صنف تكافؤ~$[w]$؛
  \item مجموعة أصناف التكافؤ تقسّم~$W$ إلى مجموعات جزئية متنافية
    اثنتين اثنتين ومستنفدة لها مجتمعة.
  \end{enumerate}
\end{prop}

\begin{prop}\ollabel{prop:S5=univ}
  تكون !!^a{formula}~$!A$ صالحة في جميع الأطر
  $\mModel{F} =\tuple{W,R}$ التي تكون~$R$ فيها علاقة تكافؤ، إذا وفقط
  إذا كانت صالحة في جميع الأطر~$\mModel{F} =\tuple{W,R}$ التي تكون
  $R$ فيها كلية. ومن ثم فإن منطق الأطر الكلية ليس إلا~$\Log{S5}$.
\end{prop}

\begin{proof}
  يسهل مباشرة التحقق من أن العلاقة الكلية~$R$ على~$W$ علاقة تكافؤ.
  ومن ثم، إذا كانت~$!A$ صالحة في جميع الأطر التي تكون~$R$ فيها علاقة
  تكافؤ، فهي صالحة في جميع الأطر الكلية. وللاتجاه الآخر نستدل
  بالمعاكس النقيض: افترض أن~$!B$ !!a{formula} تفشل في عالم~$w$ من
  نموذج~$\mModel{M} = \tuple{W, R, V}$ مبني على إطار~$\tuple{W,R}$
  تكون~$R$ عليه علاقة تكافؤ على~$W$. إذن~$\mSat/{M}{!B}[w]$. عرّف
  نموذجًا~$\mModel{M}' = \tuple{W', R', V'}$ كما يأتي:
  \begin{enumerate}
  \item $W' = [w]$؛
  \item $R'$ كلية على~$W'$؛
  \item $V'(p) = V(p) \cap W'$.
  \end{enumerate}
  (إذن تمثل المنطقة المظللة في~\olref{fig:partition} مجموعة
  العوالم~$W'$ في~$\mModel{M}'$.) ويسهل أن نرى أن~$R$ و$R'$ تتفقان
  على~$W'$. ويمكن عندئذ أن نثبت بالاستقراء على ال!!{formula}s أنه لكل
  $w' \in W'$: يصدق~$\mSat{M'}{!A}[w']$ إذا وفقط إذا كان
  $\mSat{M}{!A}[w']$، لكل~$!A$ (وهذا ذو معنى لأن~$W' \subseteq W$).
  وبخاصة،~$\mSat/{M'}{!B}[w]$، فتفشل~$!B$ في نموذج مبني على إطار كلي.
\end{proof}

\begin{figure}[t]
  \centering
  \begin{tikzpicture}[node distance=2cm, auto, thick]
    \clip [rounded corners] (0,0) -- (6,0) -- (6,4) -- (0,4) --  cycle;
    \begin{scope}
      \clip (0,0) -- (2,0)  .. controls (1.5,1.5) and (4.5,1.5)
      .. (4,4) -- (0,4) -- (0,0) -- cycle;
      \filldraw[fill=gray!40] (0,4) -- (0,2) .. controls (1.5,1.5)
      and (4.5,1.5) .. (4.5,0) -- (6,0) -- (6,4) -- (0,4) --  cycle;
    \end{scope}
    \draw [name path=line1] (2,0) .. controls (1.5,1.5) and (4.5,1.5) .. (4,4);
    \draw [name path=line2] (0,2)  .. controls (1.5,1.5) and (4.5,1.5) .. (4.5,0);
    \path [name intersections={of=line1 and line2}];
    \draw [rounded corners, use as bounding box, very thick] (0,0)
    -- (6,0) -- (6,4) -- (0,4) --  cycle;
    \node at (2,3) {$[w]$} ;
    \node at (1,1) {$[u]$} ;
    \node at (3,0.75) {$[v]$} ;
    \node at (4.75,2) {$[z]$} ;
  \end{tikzpicture}
  \caption{تقسيم~$W$ إلى أصناف تكافؤ.}
  \ollabel{fig:partition}
\end{figure}

\end{document}
